% Figure | used in parts/01-algebra/07-rotors/01-reflections-and-rotors.tex
%!TEX root = ../main.tex
% REWRITTEN: rotor R -> \rotor in the caption (docs/NOTATION.md)
\begin{figure}[ht]
\centering
\begin{tikzpicture}[>=Stealth,scale=1.1]
  \draw[thick,blue!60!black] (190:3.2)--(10:3.2) node[right]{mirror $1$ (normal $a$)};
  \draw[thick,orange!80!black] (220:3.2)--(40:3.2) node[right]{mirror $2$ (normal $b$)};
  \draw[->,thick] (0,0)--(60:2.4) node[above]{$x$};
  \draw[->,thick,blue!60!black] (0,0)--(-40:2.4) node[below right]{$x'=-axa^{-1}$};
  \draw[->,thick,red] (0,0)--(120:2.4) node[above left]{$x''=(ba)\,x\,(ba)^{-1}$};
  \draw[red,->] (60:0.9) arc[start angle=60,end angle=120,radius=0.9];
  \node[red] at (90:1.25) {$2\theta$};
  \draw[orange!80!black,->] (10:1.5) arc[start angle=10,end angle=40,radius=1.5];
  \node[orange!80!black] at (25:1.8) {$\theta$};
\end{tikzpicture}
\caption{Two reflections compose to a rotation by \emph{twice} the angle $\theta$ between the mirrors. The rotor is $\rotor=ba=e^{-B\theta}$.}
\label{fig:tworefl}
\end{figure}
